\begin{conf-abstract}
{Dissolved Gas Monitoring at the Bedretto Underground Laboratory: Responses to Seismic Stimulation}
{Alexandra Lightfoot, Matthias Brennwald, Valentin Gischig, Alba Zappone, Rolf Kipfer, Stefan Wiemer}
{SED ETH Zürich}
{Gas Monitoring at the Bedretto Underground Laboratory: Responses to Seismic Stimulation In hydraulic stimulation experiments, such as those used in seismic and geothermal energy studies, monitoring dissolved reactive and noble gases in groundwater can provide key insights into fluid transport and redistribution. While tracers added to injection water reveal specific flow paths, dissolved gas monitoring in general captures a broader picture of how fluids are mobilised or redirected under pressure. At the Bedretto Underground Laboratory for Geosciences and Geoenergy (BULGG), we tested the geochemical response during an experiment designed to trigger a controlled magnitude-zero earthquake. High-pressure water (up to 20 MPa) was injected into borehole ST1, with krypton added as a tracer. A nearby borehole (ST2), 44 m away and discharging continuously, was equipped with a miniRuedi to measure dissolved He, Ar, Kr, N2, O2, and CO2 in real time, complemented by parallel radon monitoring. This setup enabled us to track both immediate gas responses to pressure and delayed tracer migration. Dissolved gas signals at ST2 revealed a clear response once injection began. N2 and Ar increased in terms of partial pressure, while He decreased, showing that injection resulted in mobilised or redirected younger, less radiogenic fluids toward ST2 well before the Kr tracer arrived. Radon activity increased only after tracer breakthrough, consistent with fracture opening and fluid movement under stress. These findings highlight dissolved gas monitoring as a sensitive tool for probing the interplay between stress, fractures, and fluids, particularly in the context of hydraulic stimulation.}
\end{conf-abstract}
